\documentclass[11pt]{article}
\usepackage[T1]{fontenc}
\usepackage[utf8]{inputenc}
\usepackage{lmodern,amsmath,amssymb,amsthm,mathtools}
\usepackage[margin=1in]{geometry}
\usepackage{microtype,enumitem,booktabs,tabularx}
\usepackage[hidelinks]{hyperref}
\newtheorem{theorem}{Theorem}[section]
\newtheorem{proposition}[theorem]{Proposition}
\newtheorem{lemma}[theorem]{Lemma}
\theoremstyle{definition}
\newtheorem{definition}[theorem]{Definition}
\newtheorem{remark}[theorem]{Remark}
\newcommand{\R}{\mathbb R}
\newcommand{\Pow}{\mathcal P}
\newcommand{\id}{\operatorname{id}}
\newcommand{\conv}{\operatorname{conv}}
\newcommand{\rank}{\operatorname{rank}}
\newcommand{\Mod}{\operatorname{Mod}}
\newcommand{\Sen}{\operatorname{Sen}}
\newcommand{\Th}{\operatorname{Th}}
\newcommand{\Obs}{\operatorname{Obs}}
\newcommand{\eps}{\varepsilon}
\newcommand{\one}{\mathbf 1}
\setlist{nosep,leftmargin=*}
\setlength{\parskip}{3pt}
\title{Formulations of categorical phenomenology compared through a noisy XOR experiment}
\author{Technical comparison note}
\date{4 November 2026}
\begin{document}
\maketitle
\begin{abstract}
The choice of a formulation changes what is supplied, what is sought, and which proof would answer the question. A classical four-state experiment makes this distinction explicit. Two noisy marginal tests generate a square of response profiles. Adding one controlled context produces a tetrahedron, restores an exact action, and repairs a failed comparison between native composition and composition through a coarse boundary. The same example supports five further formulations: justified theories, sound abstract knowledge transformers, predictive coalgebras, categories of compatible realizations, and comparison of statistical experiments. These formulations are compatible but ask for different outputs. The strongest additional perspective for this example is sound abstraction: an honest coarse prediction exists even when exact coarse dynamics does not, and a precise completeness condition identifies the difference. All displayed results are elementary derivations or applications of established mathematics. No novelty claim or exhaustive foundational classification is made.
\end{abstract}

\section{Purpose and relation to the existing programme}
An observation-to-structure construction and a two-sided context-sensitive realization need not compete. A reconstruction can be defined from a preparation--test pairing, and its universal property can be a two-sided factorization property. The distinction becomes consequential when choosing a target: a canonical response invariant, a native implementation, a predictive state, a justified statement, a coherent family, and a finite-evidence guarantee are different outputs.

Category theory does not select factorization systems uniquely. Adjunctions, logical satisfaction, final coalgebras, categories over a base, and orders of experimental informativeness are also categorical structures. Conversely, beginning with experience does not determine the available experimental language, admissible model class, algebraic operations, or interpretation of uncertainty.

The existing programme already contains more than the initial two descriptions. Its generated-image results, predictive filtering and minimization, positive Hankel and cone realizations, native stochastic constraints, coherent cylinder comparisons, joint realization results, and monodromy examples should be preserved. In particular, the comparison below does not present generic observation-change coherence, an image theorem, or a rank/nonnegative-rank gap as an unoccupied research problem. It asks which existing mechanisms should be used for which question.

\section{The experiment and its access conditions}
\label{sec:setup}
Let
\[
 \Lambda=\{00,01,10,11\},\qquad S=\Delta\Lambda,
 \qquad \eps=\tfrac1{10},\qquad \lambda=1-2\eps=\tfrac45.
\]
All convex mixtures of the four named preparation procedures are stipulated to be available. A distribution $\mu\in S$ describes a preparation law over the four states. This concrete description specifies the example; the reconstruction is not a claim to have established a true latent ontology.

A \emph{base trial} draws a fresh sample from $\mu$, chooses exactly one terminal sensor, and records its single binary result. Sensor $X$ reads the first bit with flip probability $\eps$; sensor $Y$ reads the second bit with the same flip probability. Separate trials are independent conditional on the requested preparation and test. Base data contain only ungated trials. There is no joint or sequential $X,Y$ observation of the same held sample in base access. Noise dependence between different sensors on one sample is therefore not specified or needed.

Supply the controlled gate
\[
 g(i,j)=(i\mathbin{\oplus}j,j),\qquad F=g_*:S\longrightarrow S,
 \qquad F^2=\id_S.
\]
It is applied before the one terminal sensor on that trial. The enlarged interface admits the test $Z$: apply $g$, then read $X$. Write
\[
 a=\mathbb E_\mu[i],\qquad b=\mathbb E_\mu[j],\qquad
 c=\mathbb E_\mu[i\mathbin{\oplus}j],
\]
so the exact response maps are
\begin{equation}\label{eq:profiles}
 q(\mu)=(x,y)=(\eps+\lambda a,\eps+\lambda b),\qquad
 q^+(\mu)=(x,y,z)=(x,y,\eps+\lambda c).
\end{equation}
Arbitrary finite strings of the gate are allowed before the terminal sensor. The identity $g^2=\id$ reduces these tests to $X,Y,Z$.

\paragraph{What is assumed and what is observed.}
The calibrated sensor laws, mixture preservation, gate involutivity, and protocol restrictions are exact hypotheses. Finite noisy trials estimate probabilities; they do not by themselves establish these equations for all preparations and all continuations. An alternative empirical presentation may supply the complete exact contextual profiles and context identities directly, without a latent-bit presentation. It must still specify its operations and their laws.

\paragraph{Three different kinds of uncertainty.}
A point $\mu\in S$ is a probability law for physical outcomes. A subset $K\subseteq S$ can represent epistemic uncertainty about which law applies. A probability measure on $S$, used as a prior or posterior, is a third object. They have different types. Neither a convex mixture of preparation laws nor a square of response profiles is automatically an epistemic set or a Bayesian posterior. In the binary testing problem below, the preparation source is unknown; a known label revealing which source was used is not supplied as extra evidence.

\section{Exact reconstruction and native realization}
\subsection{Square and tetrahedron}
The four pure $q$-profiles are the corners of
\[
 C=[\eps,1-\eps]^2.
\]
In the order $00,01,10,11$, their $q^+$-profiles are
\[
 (\eps,\eps,\eps),\quad (\eps,1-\eps,1-\eps),\quad
 (1-\eps,\eps,1-\eps),\quad (1-\eps,1-\eps,\eps).
\]
Let $T$ be their convex hull. The response restriction $\pi:T\to C$ forgets $z$ and satisfies $\pi q^+=q$.

\begin{proposition}\label{prop:inverse}
The map $q:S\to C$ is surjective. The map $q^+:S\to T$ is an affine isomorphism, with inverse
\begin{align*}
 \mu_{00}&=1-\frac{a+b+c}{2},&
 \mu_{01}&=\frac{-a+b+c}{2},\\
 \mu_{10}&=\frac{a-b+c}{2},&
 \mu_{11}&=\frac{a+b-c}{2},
\end{align*}
where $a=(x-\eps)/\lambda$, $b=(y-\eps)/\lambda$, and $c=(z-\eps)/\lambda$. The normalized profile matrices with columns $1,x,y$ and $1,x,y,z$ have ranks three and four respectively.
\end{proposition}
\begin{proof}
Every pair of marginal means in $[0,1]^2$ has a joint distribution, for example the independent product law. The inverse formulas follow by solving
$a=\mu_{10}+\mu_{11}$, $b=\mu_{01}+\mu_{11}$, $c=\mu_{01}+\mu_{10}$ and normalization. They are nonnegative exactly on $T$. The four fine profiles are therefore affinely independent. The square has affine dimension two. Adding the constant coordinate raises these affine dimensions by one, giving the stated ranks.
\end{proof}
Define
\[
 d=\tfrac12(\delta_{00}+\delta_{11}),\qquad
 e=\tfrac12(\delta_{01}+\delta_{10}),\qquad
 u=\tfrac14\sum_{s\in\Lambda}\delta_s.
\]
Then
\begin{equation}\label{eq:de}
 q(d)=q(e)=q(u)=(\tfrac12,\tfrac12),\qquad z(d)=\eps,
 \qquad z(e)=1-\eps.
\end{equation}
Every pure preparation was already distinguished by $q$. The newly exposed distinction is between available mixtures.

Under the specified convex doctrine, the surjection $S\to C$ followed by the inclusion $C\hookrightarrow[0,1]^2$ is a valid generated-observable factorization. The fine version returns $T$. This is a genuine common reconstruction, not merely an analogy. It does not assert closure under $F$ or membership in the class of finite free simplices.

\subsection{A strictly positive width calculation}
Select the base sensor uniformly and include its identity in the outcome. This encodes two alternative tests as one stochastic table; it does not perform two readings on one sample. The response row is
\[
 H_\mu=\tfrac12(1-x,x,1-y,y).
\]
For the four preparations,
\begin{equation}\label{eq:H}
 H=\tfrac1{20}
 \begin{pmatrix}9&1&9&1\\9&1&1&9\\1&9&9&1\\1&9&1&9\end{pmatrix}.
\end{equation}
A finite native realization of this experiment is $H=AB$, with $A$ and $B$ row-stochastic; its width is the number of rows of $B$.

\begin{proposition}\label{prop:width}
The matrix $H$ has ordinary rank three and minimum stochastic width four. Its nonnegative rank is also four. If the three fine tests are selected uniformly, the resulting table $H^+$ has ordinary rank and minimum stochastic width both four.
\end{proposition}
\begin{proof}
The row map $(x,y)\mapsto H_\mu$ is affine and injective, so $H$ has rank three. The factorization $H=I_4H$ gives width at most four. Suppose a width-three factorization existed. The rows of $H$ would lie in the triangle spanned by the three rows of $B$. Their affine dimension is two, so $\operatorname{aff}(B)=\operatorname{aff}(H)$. This plane intersects the output simplex in the square of probability rows whose two sensor blocks each sum to $1/2$. Under its affine coordinates $(x,y)$, the triangle lies in $[0,1]^2$ and contains $[.1,.9]^2$. The latter has area $16/25$, whereas a triangle contained in the unit square has area at most $1/2$. For the last assertion, the absolute determinant formula for area is convex in each vertex separately, so its maximum is attained with vertices at square corners.

Any nonnegative factorization of a stochastic matrix can be made row-stochastic without increasing width: remove zero rows of $B$, divide each remaining row by its positive sum, and multiply the corresponding column of $A$ by that sum. The resulting $A$ has row sums one because $AB=H$. Thus the same minimum is the nonnegative rank. Finally
$H^+_\mu=\tfrac13(1-x,x,1-y,y,1-z,z)$ has rank four by Proposition~\ref{prop:inverse}, while the four-state realization supplies its upper bound.
\end{proof}
The area argument is a sufficient small-noise certificate: it works whenever $(1-2\eps)^2>1/2$. It does not establish width four for every interior square. Nonnegative rank and its polytope interpretation are standard~\cite{CohenRothblum}.

The coarse square cannot be both the algebraically reachable and observable image of a finite free simplex: such a realization would be an affine isomorphism from a simplex onto a square. This does not exclude the width-four native realization just proved. Coarse and fine native minima agree, although their generated convex carriers differ.

\section{A legitimate contextual comparison}
\label{sec:context}
\begin{proposition}\label{prop:noaction}
No function $G:C\to C$ satisfies $Gq=qF$. The fine carrier has an affine involution
\[
 G^+:T\longrightarrow T,\qquad G^+(x,y,z)=(z,y,x),
 \qquad G^+q^+=q^+F.
\]
\end{proposition}
\begin{proof}
Equation~\eqref{eq:de} gives equal coarse inputs but
$q(Fd)=(\eps,1/2)$ and $q(Fe)=(1-\eps,1/2)$. Thus even set-theoretic descent fails. The fine formula follows from the test identities $Xg=Z$, $Zg=X$, and $Yg=Y$.
\end{proof}
This failure does not refute a two-sided formulation. Such a formulation may retain a relation, a span, or the actual implementation witnesses rather than demand a function. Define
\[
 R_q=\{(q\mu,qF\mu):\mu\in S\}\subseteq C\times C,
 \qquad C\xleftarrow{q}S\xrightarrow{qF}C.
\]
The relation is the supported image of the displayed span. In rescaled coordinates it has the following exact description.

\begin{proposition}\label{prop:relation}
One has $(a,b)\,R_q\,(c,b)$ if and only if
\begin{equation}\label{eq:frechet}
 |a-b|\leq c\leq\min(a+b,2-a-b).
\end{equation}
Although $F^2=\id_S$, $R_q\circ R_q$ strictly contains the identity relation of $C$.
\end{proposition}
\begin{proof}
Put $t=\mu_{11}$. The joint probabilities are
$(1-a-b+t,b-t,a-t,t)$, so feasibility is
$\max(0,a+b-1)\leq t\leq\min(a,b)$. Substituting $c=a+b-2t$ gives~\eqref{eq:frechet}; reversing the substitution proves sufficiency. Every $q\mu$ has a two-step self-return witnessed by $\mu,F\mu$. For strictness, $Fu=u$ and $q(u)=q(d)$ permit a first step witnessed by $u$ and a second by $d$. The result is the pair $((1/2,1/2),(\eps,1/2))$, which is not an identity pair.
\end{proof}
The exact failed comparison is visible before quotienting witnesses. Define
\[
 B_q=\{(\mu,\nu)\in S\times S:qF\mu=q\nu\},\qquad
 j_q:S\longrightarrow B_q,\quad j_q(\mu)=(\mu,F\mu).
\]
The map $j_q$ is injective. It respects the starting and ending coarse boundaries because $F^2=\id$. It is not surjective: $(u,d)\in B_q$ but $d\ne Fu$. Native sequential composition requires the actual intermediate equality $\nu=F\mu$; coarse composition requires only its image under $q$. Replacing $q$ by $q^+$ makes $j_{q^+}$ bijective, since $q^+$ is injective.

Thus using spans retains a useful typed comparison and its explicit obstruction, but composing over the wrong coarse boundary still permits the false splice. One may retain the native boundary, prove the relevant exactness comparison, or intentionally accept a sound uncertainty semantics. None is supplied by renaming a relation as a realization. Cylinder factorization and equipment formalisms provide substantial settings for expressing comparisons~\cite{Garner,CruttwellShulman}; the particular preservation claim still needs its own hypotheses.

\subsection{Marginal products do not select joint mechanisms}
The square $C$ is the categorical product of two intervals in convex spaces. The physical preparation simplex $S=\Delta(\{0,1\}\times\{0,1\})$ also retains correlations and is not that product of marginal state spaces.

\begin{proposition}\label{prop:section}
The marginal map $q:S\to C$ has no affine section.
\end{proposition}
\begin{proof}
Each square corner has exactly one lift, the corresponding point mass. An affine section would send the square center, written as either diagonal midpoint, to both $d$ and $e$. These are distinct.
\end{proof}
Independence supplies a continuous nonlinear section with weights
$((1-a)(1-b),(1-a)b,a(1-b),ab)$. Its plug-in action is
$(a,b)\mapsto(a+b-2ab,b)$. At $(a,b)=(1/5,3/10)$, two applications send the first coordinate to $19/50$ and then $113/250$, not $1/5$. The section silently chooses a new coupling between steps. A product decomposition of observable coordinates is therefore a valid finding, but it is not an inferred independent mechanism.

\section{Five typed alternatives}
\label{sec:five}
The following formulations can coexist with the generated images and contextual spans above. They are distinguished by outputs and proof obligations rather than by mutually exclusive philosophical labels.

\subsection{Observations to justified theories}
\textbf{Input.} A signature $\Sigma$, an explicitly specified class $\Mod(\Sigma)$ of admissible models, a satisfaction relation, and sentences $T_d$ recording exact data. \textbf{Output.} The compatible models $\Mod(\Sigma,T_d)$ and their common consequences $\Th(\Mod(\Sigma,T_d))$. The task can return justified claims without selecting one model.

In institution notation~\cite{GoguenBurstall}, signatures form a category, $\Sen$ is covariant and $\Mod$ contravariant. A signature morphism $\sigma:\Sigma\to\Sigma'$ satisfies
\[
 M'\models_{\Sigma'}\Sen(\sigma)(\varphi)
 \quad\Longleftrightarrow\quad
 \Mod(\sigma)(M')\models_\Sigma\varphi.
\]
This is a disciplined vocabulary-change condition, not an automatic completeness or inference algorithm. The desired proof may instead concern soundness, expressive adequacy, conservative extension, or effective elimination.

\textbf{XOR output.} Given the two means, the strongest parity interval is~\eqref{eq:frechet}. At the center, all $c\in[0,1]$ remain possible. The exact $Z$ law fixes the missing parameter. A chosen independence assumption must appear in $T_d$ as extra information. The interval calculation is ordinary linear elimination, not a new institution theorem. Keeping all models makes the prior model class explicit; it does not remove that assumption.

\subsection{Observations to sound knowledge transformers}
\textbf{Input.} A concrete space $S$, a surjective observable map $q:S\to C$, and an action $F:S\to S$. The epistemic domains used here are powersets, ordered by inclusion:
\[
 \alpha:\Pow(S)\to\Pow(C),\quad \alpha(K)=q[K],\qquad
 \gamma:\Pow(C)\to\Pow(S),\quad\gamma(A)=q^{-1}[A].
\]
They satisfy $\alpha\dashv\gamma$, $\alpha\gamma=\id$, and $K\subseteq\gamma\alpha K$. Let $F_!$ denote direct image on subsets.

\textbf{Output.} The best sound abstract transformer
\begin{equation}\label{eq:best}
 F^\#=\alpha F_!\gamma:\Pow(C)\longrightarrow\Pow(C).
\end{equation}
A transformer $H$ is sound if $\alpha F_!\subseteq H\alpha$ pointwise. Substituting $K=\gamma A$ shows $F^\#\subseteq H$. Hence~\eqref{eq:best} is optimal among sound transformers on this particular domain. Its existence does not imply an effective finite representation.

\textbf{Proof obligations.} Soundness, best precision, completeness $\alpha F_!=F^\#\alpha$, and exact composition are distinct. The domain may be refined, the operations restricted, or sound imprecision accepted. These are standard abstract-interpretation questions~\cite{CousotCousot,GiacobazziRanzatoScozzari}.

\textbf{XOR output.} On singleton coarse knowledge, $F^\#\{v\}=R_q[v]$. The two-step knowledge transformer strictly overapproximates the identity. This is epistemic loss from reabstracting, not extra physical randomness and not a failure of soundness. With $q^+$, no such loss occurs. Proposition~\ref{prop:bridge} proves the exact completeness criterion without assuming a desired quotient action exists.

\subsection{Histories and actions to predictive behavior}
\textbf{Input.} A Moore coalgebra $c:S\to O\times S^A$ specifying output $o:S\to O$ and transitions $F_a:S\to S$ for an action alphabet $A$. \textbf{Output.} The behavior map
\[
 \beta:S\longrightarrow O^{A^*},\qquad
 \beta(s)(w)=o(F_w(s)).
\]
This target is the final Moore coalgebra. Its image identifies precisely those states indistinguishable by every allowed future experiment. Reachability from stipulated preparations is a further restriction.

\textbf{Proof obligations.} A selected finite test vocabulary must separate the intended behavioral classes and be closed under the necessary continuations. Finite native stochastic realization is a separate issue. The explicit Moore setting avoids an unsupported claim that every stochastic-system functor has this same final object~\cite{Rutten}.

\textbf{XOR output.} Take $O=C$, $o=q$, and $A=\{g\}$. The future context $Xg$ separates $d,e$, and the complete response image is represented by $T$, with action $(x,y,z)\mapsto(z,y,x)$. Current response equality was not future-stable. This is already the kind of distinction treated by the programme's predictive/filtering work; it is an existing alternative framing, not a newly discovered missing foundation.

\subsection{Data to categories of compatible realizations}
\textbf{Input.} A specified category of native models and observation functor
$\Obs:\Mod\to\mathcal D$, with selected data $d\in\mathcal D$. \textbf{Output.} The strict fiber when equality is intended, or the isomorphism fiber with objects $(M,\phi:\Obs(M)\cong d)$ when data are considered up to isomorphism. Morphisms preserve the indicated observation identification. Its groupoid core retains only invertible comparisons; this is appropriate only when these are the intended equivalences.

\textbf{Proof obligations.} Nonemptiness, identifiability, representability, resource bounds, and a chosen universal object are different questions. Parametrized models additionally require a base/site and restriction operations before asking for effective descent or a continuous choice. Calling the assignment a stack does not establish the descent axiom~\cite{Stacks}.

\textbf{XOR output.} Within the fixed four-state ontology, the fiber over the central coarse profile is an interval of distributions. The fine data select one distribution. There is no nontrivial stack obstruction in this single finite example. The programme's monodromy examples give the relevant family-level warning: pointwise minimal realizations can fail to admit a global continuous minimal labelled choice. Constructing a fiber is organizational; proving native descent or a topological obstruction is substantive. A formal completion must not be substituted silently for the specified native model class.

\subsection{Experiments to information and decision guarantees}
\textbf{Input.} A hypothesis class $\Theta$, channels $E:\Theta\to\mathrm{Dist}(Y)$, permitted sampling protocols, and a loss or target property. \textbf{Output.} Informativeness comparisons, hypothesis tests, confidence sets, or decision guarantees. Exact garbling means $E\succeq E'$ when $E'=KE$ for a hypothesis-independent stochastic kernel $K$.

\textbf{Proof obligations.} Identifiability depends on access; quantitative guarantees depend on sampling and noise assumptions. A proposed comparison must specify its parameter space and whether it is prior-dependent. Categorical probability already supplies substantial comparison theory~\cite{FritzEtAl}; the elementary argument below does not invoke representability-dependent theorems inside $\mathbf{FinStoch}$.

\begin{proposition}\label{prop:testing}
Suppose the preparation source is one of $d,e$, with equal prior probabilities, and each permitted trial independently restarts that unknown source. Every adaptive sequence of base $X,Y$ trials has the same transcript distribution under both hypotheses. Its optimal error is $1/2$. One $Z$ trial has optimal error $\eps=1/10$.
\end{proposition}
\begin{proof}
Conditional on any preceding transcript and the chosen base sensor, the next outcome is Bernoulli$(1/2)$ under both sources. Induction proves equality of all finite transcript laws, including decisions made from them. The same conclusion holds for any well-defined stopping procedure because its stopped transcript is a function of the full identical transcript law and independent algorithmic randomness. For $Z$, the laws are Bernoulli$(\eps)$ and Bernoulli$(1-\eps)$; decide $d$ on outcome zero and $e$ on one. This has error $\eps$, and the likelihood-ratio rule shows optimality.
\end{proof}
Joint or sequential same-sample readings would change the hypothesis-testing problem. This finite comparison measures the value of an added context; it does not identify an arbitrary hidden mechanism. A confidence set $K_n\subseteq S$ with coverage guarantees is again an epistemic object, not a replacement for an underlying probability law $\mu$.

\section{A concrete bridge between the formulations}
\begin{proposition}\label{prop:bridge}
For a surjection $q:S\to C$ and function $F:S\to S$, the following are equivalent:
\begin{enumerate}
 \item There is a function $G:C\to C$ with $Gq=qF$.
 \item If $q(s)=q(t)$, then $qF(s)=qF(t)$.
 \item The relation $\{(q(s),qF(s)):s\in S\}$ is single-valued.
 \item The powerset abstraction is complete: $\alpha F_!=F^\#\alpha$.
\end{enumerate}
When $S,C$ are convex spaces and $q,F$ affine, the unique descended $G$ is affine.
\end{proposition}
\begin{proof}
The first three equivalences are the usual well-definedness condition on a quotient. For a singleton $\{s\}$, completeness says
\[
 \{qF(s)\}=\{qF(t):q(t)=q(s)\},
\]
which is precisely the second condition. Conversely this equality on every fiber gives completeness on arbitrary subsets by taking unions. In the convex case, choose preimages $s,t$ of $v,w$. Then
$G(rv+(1-r)w)=qF(rs+(1-r)t)=rG(v)+(1-r)G(w)$.
\end{proof}
The theorem concerns exact coarse dynamics; it neither identifies $S$ from observations nor supplies a native implementation of $C$.

\begin{proposition}\label{prop:testclosure}
Let $\Lambda$ be finite, $S=\Delta\Lambda$, and $W\subseteq\R^\Lambda$ a linear space of real tests containing the constant function $1$. Let $q_W(\mu)$ record the expectations of a basis of $W$, with codomain its image, and let $F$ be induced by a stochastic transition on $\Lambda$. Then the equivalent conditions of Proposition~\ref{prop:bridge} hold for $q_W$ if and only if $F^*W\subseteq W$, where
$\mathbb E_{F\mu}[f]=\mathbb E_\mu[F^*f]$.
\end{proposition}
\begin{proof}
If $F^*W\subseteq W$, equal expectations on $W$ remain equal after $F$. Conversely let $v\in W^\perp$. Since $1\in W$, the signed vector $v$ has total mass zero. Unless it is zero, it is a positive scalar multiple of a difference of two probability distributions with identical $W$ expectations. Invariance of those equivalence classes implies $\langle v,F^*f\rangle=0$ for every $f\in W$. Finite-dimensional annihilator duality gives $F^*f\in W$.
\end{proof}
For a finite supplied family of stochastic actions $F_a$, the smallest invariant test space containing $W_0$ is
\[
 W_* = \operatorname{span}\{F_w^*f:w\in A^*,\ f\in W_0\}.
\]
It can be obtained by iterating $W_{n+1}=W_n+\sum_{a\in A}F_a^*W_n$. Every strict change raises dimension, so at most $|\Lambda|-\dim W_0$ strict increases occur. This is the smallest exact refinement \emph{among expectation quotients from linear test spaces containing $W_0$}; it is not a minimum over every possible observational doctrine or experimental cost.

For the XOR experiment,
\[
 W_0=\operatorname{span}\{1,i,j\},\qquad
 g^*i=i\mathbin{\oplus}j=i+j-2ij,
\]
so
\[
 W_*=\operatorname{span}\{1,i,j,i\mathbin{\oplus}j\}=\R^\Lambda.
\]
The known affine noise transformation does not change these spans. Thus one closure step gives the tetrahedral refinement. This is a standard observability/completeness mechanism, not a new general theorem. It shows concretely how changing the framing directs proof search toward invariant test spaces rather than an impossible coarse action.

\section{What the typed questions actually request}
The following missing objects should not be interchanged:
\begin{enumerate}
 \item A response $r(p,Z)\in[0,1]$, supplied by a new measurement or a justified model calculation.
 \item A function $G:C\to C$ satisfying $Gq=qF$, whose existence is refuted here.
 \item A compatible coupling $\mu\in S$ with prescribed marginals, usually nonunique.
 \item A native compositional witness satisfying $\nu=F\mu$, stronger than the coarse equality $q\nu=qF\mu$.
 \item Stochastic matrices $A,B$ of a specified width with $H=AB$, with width three refuted and width four witnessed here.
 \item An actual implementation comparison or decomposition cell. Equal response tables do not fabricate one or equate distinct implementation witnesses without an explicit quotient.
 \item A statistical decision or confidence statement with its stated error or coverage, which cannot be replaced by an exact algebraic image theorem.
\end{enumerate}
A categorical framework can organize these jointly. Its usefulness is tested by whether it transports the intended witnesses and laws, not by whether all of them can be drawn using similar arrows. The difficult part can lie in native membership, composition, global selection, empirical identification, or computational tractability, depending on the chosen output.

\section{Controls and limits}
\begin{enumerate}
 \item \textbf{Pure-state separation.} The base tests separate every pure preparation. The obstruction requires the stipulated operational mixtures.
 \item \textbf{Action control.} Identity or exchange of the two coordinates preserves $W_0$. The same marginal abstraction then supports exact dynamics.
 \item \textbf{Access control.} Adding joint or sequential same-sample $X,Y$ observations can expose correlation without XOR. The base statistical impossibility applies only to the declared fresh-sample interface.
 \item \textbf{Noise control.} At $\eps=1/2$ all sensors are uninformative. Unknown calibration is a different identification problem. The width-four area certificate also has its explicit small-noise threshold.
 \item \textbf{Precision control.} Returning every coarse state is sound but needlessly imprecise. Best sound approximation and completeness must be assessed separately.
 \item \textbf{Coupling control.} Choosing independence is additional modeling structure. Proposition~\ref{prop:section} and the two-step calculation show what it fails to preserve.
 \item \textbf{Ontology control.} Injectivity of $q^+$ identifies the preparation law inside the stipulated four-state class, not the unique physical ontology among all implementations.
 \item \textbf{Contextuality control.} This is a fully classical experiment with an explicit global model. The word context denotes controlled preprocessing. No Bell or Kochen--Specker no-go claim is being made.
 \item \textbf{Empirical control.} Exact context laws were supplied. Finite positive-probability data alone do not prove exact closure or all-continuation equivalence.
 \item \textbf{Categorical control.} The coarse convex reconstruction and span are legitimate. Their existence should not be confused with exact preservation of a particular native operation or boundary gluing.
\end{enumerate}

\section{Assessment and source status}
The strongest additional formulation in this comparison is observation to justified claims together with sound knowledge transformers. It distinguishes an impossible exact quotient action from an available, useful uncertainty semantics and makes the repair condition precise. Predictive coalgebra provides a compatible operation-closed reconstruction. Categories of realizations retain nonuniqueness and global obstructions. Statistical comparison asks what extra experiments buy in evidence rather than in formal structure.

These descriptions do not exhaust categorical phenomenology, and none should replace the programme's existing comparison theorems by assertion. The practical choice is the intended output. For this example, a common target is to relate exact quotient dynamics, invariant observation equivalence, complete abstraction, and action-closed tests. Native realization width and statistical guarantees remain separate proved statements.

The constructions, source roles, and limits are explicit: Garner supplies cylinder-factorization theory; abstract interpretation supplies soundness and completeness; coalgebra supplies future behavior; institution theory supplies satisfaction under vocabulary change; descent requires effective gluing; categorical statistics supplies experimental comparison. The finite XOR derivations and counterexamples are proved above. The source review is bounded and gives no novelty clearance.

\begin{thebibliography}{99}
\bibitem{Garner}
R. Garner, \emph{The Isbell monad}, 2014.
\url{https://arxiv.org/abs/1410.7108}.

\bibitem{CohenRothblum}
J. E. Cohen and U. G. Rothblum,
\emph{Nonnegative Ranks, Decompositions, and Factorizations of Nonnegative Matrices},
Linear Algebra and its Applications 190 (1993), 149--168.
\url{https://lab.rockefeller.edu/cohenje/PDFs/211CohenRothblumLAA1993.pdf}.

\bibitem{CruttwellShulman}
G. S. H. Cruttwell and M. A. Shulman,
\emph{A unified framework for generalized multicategories},
Theory and Applications of Categories 24 (2010), 580--655.
\url{https://arxiv.org/abs/0907.2460}.

\bibitem{GoguenBurstall}
J. A. Goguen and R. M. Burstall,
\emph{Institutions: Abstract Model Theory for Specification and Programming},
Journal of the ACM 39 (1992), 95--146.
\url{https://doi.org/10.1145/147508.147524}.
Author bibliography and original-paper link:
\url{https://cseweb.ucsd.edu/~goguen/projs/sem.html}.

\bibitem{CousotCousot}
P. Cousot and R. Cousot,
\emph{Abstract interpretation: a unified lattice model for static analysis of programs by construction or approximation of fixpoints},
POPL 1977, 238--252.
\url{https://cs.nyu.edu/~pcousot/COUSOTpapers/POPL77.shtml}.

\bibitem{GiacobazziRanzatoScozzari}
R. Giacobazzi, F. Ranzato and F. Scozzari,
\emph{Making Abstract Interpretations Complete},
Journal of the ACM 47 (2000), 361--416.
\url{https://www.sci.unich.it/~scozzari/paper/JACM00.pdf}.

\bibitem{Rutten}
J. J. M. M. Rutten, \emph{Universal coalgebra: a theory of systems},
Theoretical Computer Science 249 (2000), 3--80.
\url{https://ir.cwi.nl/pub/48/}.

\bibitem{Stacks}
The Stacks Project Authors,
\emph{Descent data in fibred categories} and \emph{Stacks in groupoids},
Tags 02ZC and 02ZI.
\url{https://stacks.math.columbia.edu/tag/02ZC};
\url{https://stacks.math.columbia.edu/tag/02ZI}.

\bibitem{FritzEtAl}
T. Fritz, T. Gonda, P. Perrone and E. F. Rischel,
\emph{Representable Markov Categories and Comparison of Statistical Experiments in Categorical Probability},
Theoretical Computer Science 961 (2023), 113896.
\url{https://arxiv.org/abs/2010.07416}.
\end{thebibliography}
\end{document}
